\documentclass[12pt,a4paper]{article}
\usepackage[margin=2.5cm]{geometry}
\usepackage{amsmath,amssymb,amsthm}
\usepackage{booktabs,array}
\usepackage{hyperref}
\usepackage{tcolorbox}

\newtheorem{theorem}{Theorem}
\newtheorem{lemma}{Lemma}
\newtheorem{proposition}{Proposition}
\newtheorem{corollary}{Corollary}
\theoremstyle{definition}
\newtheorem{definition}{Definition}
\newtheorem{remark}{Remark}

\title{Verification or Falsification of the Claim\\
$\frac{35}{18} = \frac{7 \times 5}{2 \times 9}$ from the
Gilkey $a_4$ Gauge Prefactor\\[6pt]
\large Independent Derivation from the Seeley--DeWitt Heat Kernel}

\author{Critical Analysis for G.\ T.\ Alcock\\
Density Field Dynamics Research Campaign}
\date{March 22, 2026}

\begin{document}
\maketitle

\begin{abstract}
We investigate whether the prefactor $35/18 = 7 \times 5/(2 \times 9)$
in the DFD formula
$\kappa_{U(1)}/\beta_{U(1)} \approx (35/18)(64/65)$
can be \emph{derived} from the Gilkey $a_4$ Seeley--DeWitt coefficient
for the connection Laplacian on the bundle
$\mathcal{O}(9) \oplus \mathcal{O}^{\oplus 5}$ over $\mathbb{C}P^2$,
or whether it is reverse-engineered.  Starting from the standard Gilkey
formula (Vassilevich 2003, Gilkey 1995), we compute the gauge-kinetic
coefficient step by step, carefully distinguishing the background
curvature invariants from the gauge fluctuation terms.

\medskip\noindent
\textbf{Verdict:} The number $35/18$ does \emph{not} emerge directly
from the $a_4$ Gilkey formula applied to the connection Laplacian on
$\mathcal{O}(9) \oplus \mathcal{O}^{\oplus 5}$ over $\mathbb{C}P^2$.
The Gilkey gauge-kinetic term produces the coefficient $1/96$ for
$\mathcal{O}(1)$, which scales to $81/96 = 27/32$ for $\mathcal{O}(9)$
(with zero contribution from $\mathcal{O}^{\oplus 5}$).  No combination
of these Gilkey outputs yields $35/18$.  The factor $35/18$ is an
\emph{approximate numerical fit} that decomposes into SM quantum numbers
($N_{\rm sp} = 7$, $n = 5$, $a = 9$), but this decomposition is
\emph{a posteriori}---it was constructed to match the numerical ratio
$\kappa/\beta \approx 1.9147$, not derived from the heat kernel.
\end{abstract}

\tableofcontents
\newpage

%=============================================================================
\section{The Claim Under Scrutiny}
\label{sec:claim}
%=============================================================================

The DFD derivation chain for $\alpha^{-1}$ includes the formula:
\begin{equation}\label{eq:claim}
\frac{\kappa_{U(1)}}{\beta_{U(1)}}
\approx \frac{35}{18} \times \frac{64}{65}
= \frac{224}{117} = 1.91453\ldots
\end{equation}
where $35/18 = N_{\rm sp} \cdot n / (2a)$ with $N_{\rm sp} = 7$,
$n = 5$, $a = 9$.

The existing documents (\texttt{a4\_Gilkey\_Computation.tex},
\texttt{Kappa\_Beta\_Ratio.tex}) present this factor as arising from
``the gauge-kinetic coefficient from the $a_4$ Seeley--DeWitt formula.''

\textbf{The critique:} The ratio $\kappa/\beta = 1.91469\ldots$ is
known numerically, and $35/18 \times 64/65 = 1.91453$ matches it
to 85 ppm.  But \emph{closeness does not prove derivation}.  We must
check whether the Gilkey $a_4$ formula actually produces 35/18 when
applied to the relevant bundle, or whether this is numerology.

%=============================================================================
\section{The Gilkey $a_4$ Formula: Precise Statement}
\label{sec:gilkey}
%=============================================================================

\subsection{General Formula}

For a second-order operator $P = -(g^{ij}\nabla_i\nabla_j + E)$ on a
vector bundle $V$ of rank $r$ over a compact Riemannian manifold $M^m$
without boundary (Vassilevich, \textit{Phys.\ Rep.}\ 388 (2003) 279,
Eq.~4.1; Gilkey, \textit{Invariance Theory}, Ch.~4):
\begin{equation}\label{eq:a4-general}
\boxed{
a_4(P) = \frac{(4\pi)^{-m/2}}{360}\int_M \mathrm{tr}_V\!\left[
60\,ER + 180\,E^2 + 30\,\Omega_{ij}\Omega^{ij}
+ \bigl(12\nabla^2 R - 5R^2 + 2|{\rm Ric}|^2 + 2|{\rm Riem}|^2\bigr)
\cdot\mathrm{Id}_V
\right]d\mathrm{vol}
}
\end{equation}
where:
\begin{itemize}
\item $\nabla$ is the total connection (Levi-Civita $+$ bundle connection),
\item $E$ is the endomorphism-valued potential (zero for the connection Laplacian),
\item $\Omega_{ij} = [\nabla_i, \nabla_j] - \nabla_{[e_i, e_j]}$ is
  the bundle curvature 2-form (valued in $\mathrm{End}(V)$),
\item The trace $\mathrm{tr}_V$ runs over the fibre of $V$.
\end{itemize}

\subsection{Specialization: Connection Laplacian ($E = 0$)}

For the \textbf{connection Laplacian} $P = -g^{ij}\nabla_i\nabla_j$
(i.e., $E = 0$), the $E$-dependent terms vanish:
\begin{equation}\label{eq:a4-conn-lap}
a_4(P) = \frac{(4\pi)^{-m/2}}{360}\int_M \left[
30\,\mathrm{tr}_V(\Omega_{ij}\Omega^{ij})
+ r\bigl(12\nabla^2 R - 5R^2 + 2|{\rm Ric}|^2 + 2|{\rm Riem}|^2\bigr)
\right]d\mathrm{vol}
\end{equation}
where $r = \mathrm{rank}(V)$.

\subsection{The Gauge-Kinetic Piece}

In the spectral action approach (Chamseddine--Connes 1997), one
expands $\mathrm{Tr}\,f(D^2/\Lambda^2)$ and the gauge coupling
constant is read off from the $a_4$ coefficient.  Specifically:
\begin{equation}\label{eq:gauge-kinetic-def}
\text{gauge kinetic coefficient} = -\frac{(4\pi)^{-m/2}}{360}
\times 30 \times \int_M \mathrm{tr}_V(\Omega_{ij}\Omega^{ij})\,d\mathrm{vol}
\end{equation}

The key point: \textbf{only the $\mathrm{tr}(\Omega^2)$ term contributes
to the gauge coupling.}  The curvature invariants
$R^2, |{\rm Ric}|^2, |{\rm Riem}|^2$ contribute to the
cosmological constant and Einstein--Hilbert action, not to the
gauge kinetic term.

%=============================================================================
\section{Computation on $\mathcal{O}(1)$ over $\mathbb{C}P^2$}
\label{sec:O1}
%=============================================================================

\subsection{Bundle Curvature of $\mathcal{O}(1)$}

The tautological line bundle $\mathcal{O}(1)$ on $\mathbb{C}P^2$ has
curvature (in skew-Hermitian convention):
\begin{equation}
\Omega^{\mathcal{O}(1)} = i\omega_{\rm FS}
\end{equation}
where $\omega_{\rm FS}$ is the Fubini--Study K\"ahler form (normalized
so the holomorphic sectional curvature is $c = 4$).

Since $\mathcal{O}(1)$ is a \textbf{line bundle} ($r = 1$), the trace
is trivial:
\begin{equation}
\mathrm{tr}(\Omega_{ij}\Omega^{ij})
= \Omega_{ij}\Omega^{ij}
= (i\omega_{ij})(i\omega^{ij})
= -\omega_{ij}\omega^{ij}
\end{equation}

\subsection{Computing $|\omega|^2$ on $\mathbb{C}P^2$}

For the Fubini--Study metric on $\mathbb{C}P^n$ with holomorphic
sectional curvature $c$, the K\"ahler form satisfies:
\begin{equation}
\omega_{ij}\omega^{ij} = 2n
\end{equation}
This is because $\omega_{ij} = g_{ik}J^k{}_j$ where $J$ is the complex
structure satisfying $J^2 = -\mathrm{Id}$, so:
\begin{equation}
\omega_{ij}\omega^{ij} = g_{ik}J^k{}_j g^{ip}g^{jq}\omega_{pq}
= J^k{}_j J^j{}_k = -\mathrm{tr}(J^2) = \mathrm{tr}(\mathrm{Id}_{2n})
= 2n
\end{equation}

For $n = 2$ (i.e., $\mathbb{C}P^2$):
\begin{equation}
|\omega|^2 = \omega_{ij}\omega^{ij} = 4
\end{equation}

Therefore:
\begin{equation}
\mathrm{tr}(\Omega^{\mathcal{O}(1)}_{ij}\Omega^{\mathcal{O}(1)\,ij})
= -4
\end{equation}

\subsection{The $a_4$ Gauge-Kinetic Coefficient for $\mathcal{O}(1)$}

From \eqref{eq:gauge-kinetic-def} with $m = 4$,
$\mathrm{Vol}(\mathbb{C}P^2) = \pi^2/2$:
\begin{align}
a_4^{\rm gauge}(\mathcal{O}(1))
&= -\frac{1}{(4\pi)^2}\cdot\frac{30}{360}\cdot(-4)\cdot\frac{\pi^2}{2}
\nonumber\\
&= -\frac{1}{16\pi^2}\cdot\frac{1}{12}\cdot(-4)\cdot\frac{\pi^2}{2}
\nonumber\\
&= \frac{1}{16\pi^2}\cdot\frac{4}{12}\cdot\frac{\pi^2}{2}
\nonumber\\
&= \frac{1}{16\pi^2}\cdot\frac{2\pi^2}{12}
= \frac{1}{96}
\end{align}

\begin{equation}
\boxed{a_4^{\rm gauge}(\mathcal{O}(1),\,\mathbb{C}P^2) = \frac{1}{96}}
\end{equation}

(Note: the sign convention here gives a \emph{positive} gauge-kinetic
coefficient because $\mathrm{tr}(\Omega^2) < 0$ for a unitary connection,
and there is an extra minus sign in the spectral action coupling.)

%=============================================================================
\section{Computation on $\mathcal{O}(9)$ over $\mathbb{C}P^2$}
\label{sec:O9}
%=============================================================================

\subsection{Bundle Curvature of $\mathcal{O}(q)$}

For the line bundle $\mathcal{O}(q)$ on $\mathbb{C}P^n$:
\begin{equation}
\Omega^{\mathcal{O}(q)} = iq\,\omega_{\rm FS}
\end{equation}

Therefore:
\begin{equation}
\mathrm{tr}(\Omega^{\mathcal{O}(q)}_{ij}\Omega^{\mathcal{O}(q)\,ij})
= -q^2\,\omega_{ij}\omega^{ij} = -4q^2
\end{equation}

\subsection{For $q = 9$}

\begin{equation}
\mathrm{tr}(\Omega^{\mathcal{O}(9)}_{ij}\Omega^{\mathcal{O}(9)\,ij})
= -4 \times 81 = -324
\end{equation}

The gauge-kinetic coefficient:
\begin{align}
a_4^{\rm gauge}(\mathcal{O}(9))
&= -\frac{1}{16\pi^2}\cdot\frac{30}{360}\cdot(-324)\cdot\frac{\pi^2}{2}
\nonumber\\
&= \frac{1}{16\pi^2}\cdot\frac{324}{12}\cdot\frac{\pi^2}{2}
= \frac{27}{32}
\end{align}

\begin{equation}
\boxed{a_4^{\rm gauge}(\mathcal{O}(9),\,\mathbb{C}P^2)
= \frac{27}{32} = q^2 \cdot \frac{1}{96} \cdot 3 = \frac{81}{96}}
\end{equation}

Wait---let me recheck: $81/96 = 27/32$.  Yes, $81/96 = 27/32$. $\checkmark$

%=============================================================================
\section{The Direct Sum $\mathcal{O}(9) \oplus \mathcal{O}^{\oplus 5}$}
\label{sec:direct-sum}
%=============================================================================

\subsection{Heat Kernel of a Direct Sum}

For a direct sum of bundles $V = V_1 \oplus V_2$:
\begin{equation}
a_k(V_1 \oplus V_2) = a_k(V_1) + a_k(V_2)
\end{equation}

The bundle $\mathcal{O}^{\oplus 5}$ is the trivial rank-5 bundle.
Its connection is flat ($\Omega = 0$), so:
\begin{equation}
a_4^{\rm gauge}(\mathcal{O}^{\oplus 5}) = 0
\end{equation}

Therefore:
\begin{equation}
\boxed{
a_4^{\rm gauge}(\mathcal{O}(9) \oplus \mathcal{O}^{\oplus 5},\,\mathbb{C}P^2)
= \frac{27}{32} + 0 = \frac{27}{32}
}
\end{equation}

\subsection{What the Gilkey Formula Gives}

The Gilkey gauge-kinetic coefficient for the connection Laplacian on
$\mathcal{O}(9) \oplus \mathcal{O}^{\oplus 5}$ over $\mathbb{C}P^2$ is:
\begin{equation}
\frac{27}{32} = 0.84375
\end{equation}

\textbf{This is not $35/18 = 1.9\overline{4}$.}

%=============================================================================
\section{Attempting to Recover 35/18}
\label{sec:attempt}
%=============================================================================

Let us now try every reasonable way to extract $35/18$ from the
Gilkey computation.

\subsection{Attempt 1: Ratio to scalar Laplacian $a_4$}

The $a_4$ coefficient for the \emph{scalar} Laplacian on
$\mathbb{C}P^2$ (rank-1 trivial bundle, $\Omega = 0$) involves only
the curvature invariants:
\begin{align}
a_4^{\rm scal}(\mathbb{C}P^2)
&= \frac{1}{16\pi^2 \cdot 360}\bigl(-5R^2 + 2|{\rm Ric}|^2
+ 2|{\rm Riem}|^2\bigr)\cdot\frac{\pi^2}{2}
\nonumber\\
&= \frac{-5(576) + 2(144) + 2(192)}{360 \cdot 32}
= \frac{-2208}{11520} = -\frac{23}{120}
\end{align}

Ratio:
\begin{equation}
\frac{a_4^{\rm gauge}(\mathcal{O}(9))}{|a_4^{\rm scal}|}
= \frac{27/32}{23/120} = \frac{27 \times 120}{32 \times 23}
= \frac{3240}{736} = \frac{405}{92} \approx 4.402
\end{equation}

This is not $35/18$. $\times$

\subsection{Attempt 2: Normalize by $a_0$}

\begin{equation}
a_0(\mathbb{C}P^2) = \frac{\mathrm{Vol}}{(4\pi)^2} = \frac{1}{32}
\end{equation}

\begin{equation}
\frac{a_4^{\rm gauge}(\mathcal{O}(9))}{a_0}
= \frac{27/32}{1/32} = 27
\end{equation}

Not $35/18$. $\times$

\subsection{Attempt 3: Include $S^3$ via the product formula}

On the product $\mathbb{C}P^2 \times S^3$, the gauge field lives on
the $\mathbb{C}P^2$ factor.  The product heat-kernel convolution:
\begin{equation}
a_4^{\rm gauge}(\mathbb{C}P^2 \times S^3)
= a_4^{\rm gauge}(\mathbb{C}P^2) \times a_0(S^3)
\end{equation}
where:
\begin{equation}
a_0(S^3) = (4\pi)^{-3/2} \times 2\pi^2
= \frac{2\pi^2}{8\pi^{3/2}} = \frac{\sqrt{\pi}}{4}
\approx 0.4431
\end{equation}

So:
\begin{equation}
a_4^{\rm gauge}(\mathcal{O}(9),\,\mathbb{C}P^2 \times S^3)
= \frac{27}{32} \times \frac{\sqrt{\pi}}{4}
= \frac{27\sqrt{\pi}}{128} \approx 0.3738
\end{equation}

This is not $35/18$ either. $\times$

\subsection{Attempt 4: The ``gauge coupling'' in spectral action language}

In the spectral action, the gauge coupling $g^{-2}$ is extracted as:
\begin{equation}
\frac{1}{g^2} = \frac{f_2\Lambda^2}{2\pi^2} \times
\bigl(\text{spectral function}\bigr) \times
\bigl(\text{trace factors}\bigr)
\end{equation}
where $f_2 = \int_0^\infty t\,f(t)\,dt$ is the second moment of the
cutoff function.  The relationship between $a_4^{\rm gauge}$ and the
gauge coupling involves the specific cutoff profile, which is
\textbf{model-dependent} and cannot be determined from the Gilkey
formula alone.

\subsection{Attempt 5: Decompose $35/18$ and check each factor}

The claimed decomposition is:
\begin{equation}
\frac{35}{18} = \frac{2}{3} \times \frac{5}{2} \times \frac{7}{6}
\end{equation}
where:
\begin{itemize}
\item $2/3 = 4\pi/(6\pi)$: normalization relating $\alpha$ to $\kappa$.
\item $5/2 = 1 + R_W/4$: electroweak mixing with Wilson ratio $R_W = 6$.
\item $7/6 = N_{\rm sp}/(N_{\rm gen} \cdot n_2)$: spectral triple correction.
\end{itemize}

\textbf{Observation:} The factor $2/3$ is \emph{not} from the $a_4$
formula---it is from the definition of $\kappa$ vs.\ $\alpha$.  The
factor $5/2$ is \emph{not} from the $a_4$ formula---it is from
electroweak mixing.  The factor $7/6$ is \emph{not} from the $a_4$
formula---it is from the SM content of the spectral triple.

None of these three factors come from the Gilkey $a_4$ coefficient.
They are SM physics and normalization conventions.

\subsection{Attempt 6: Can $N_{\rm sp} \cdot n / (2a)$ emerge from $a_4$?}

The alternative parameterization is $35/18 = 7 \times 5/(2 \times 9)$.
In the $a_4$ computation:
\begin{itemize}
\item The number $9$ appears as $q^2 = 81$ in $|\Omega|^2 = 4q^2 = 324$,
  so $a = q^2 = 9$ enters through $a_4^{\rm gauge} \propto q^2 = a$.
\item The number $5$ appears as the count of trivial summands
  $\mathcal{O}^{\oplus 5}$, but these have $\Omega = 0$ and contribute
  \textbf{zero} to $a_4^{\rm gauge}$.
\item The number $7 = N_{\rm sp}$ does not appear anywhere in the
  Gilkey formula for the connection Laplacian on a fixed bundle.
  It is a Standard Model input.
\end{itemize}

\begin{tcolorbox}[colback=red!5!white, colframe=red!60!black,
  title=Key Finding]
The number $n = 5$ (trivial bundles) contributes \textbf{zero} to the
gauge-kinetic term of $a_4$, because trivial bundles have flat
connections ($\Omega = 0$).  Yet the formula $35/18 = 7 \times 5/(2 \times 9)$
\emph{requires} $n = 5$ in the numerator.

This is a mathematical impossibility: the Gilkey gauge-kinetic coefficient
$a_4^{\rm gauge}$ depends on $|\Omega|^2$, and the five trivial bundles
have $\Omega = 0$.  Therefore $n = 5$ \textbf{cannot} enter the
gauge-kinetic prefactor through the $a_4$ formula.
\end{tcolorbox}

%=============================================================================
\section{What the $a_4$ Formula Actually Produces}
\label{sec:actual}
%=============================================================================

\subsection{The Honest Gauge-Kinetic Coefficient}

From the Gilkey formula, the gauge-kinetic coefficient for a U(1)
gauge field on $\mathbb{C}P^2$ is determined by $|\Omega|^2$:
\begin{equation}
\kappa_{\rm gauge}^{\mathcal{O}(q)} = \frac{q^2}{96}
\end{equation}

For the physical U(1) gauge field (charge $q_1 = 3$, so $q = 3q_1 = 9$
for the total twist... or is $q = q_1 = 3$?  Or is $q = 1$ for the
hypercharge U(1)?), the coefficient depends on which bundle carries
the physical gauge field.

\subsection{The Normalization Ambiguity}

This is where the ``derivation'' of 35/18 becomes questionable.  In
the Chamseddine--Connes spectral action, the gauge coupling is:
\begin{equation}
\frac{1}{g^2} \sim f_2 \times \mathrm{Tr}_F(Q^2) \times a_4^{\rm gauge}
\end{equation}
where $\mathrm{Tr}_F(Q^2)$ is a trace over the \emph{finite-dimensional}
internal Hilbert space $\mathcal{H}_F$ of the spectral triple.  This
trace involves the Standard Model quantum numbers and produces factors
like $N_{\rm sp}$, $\mathrm{Tr}(Y^2)$, and $g_F$.

The factor $35/18$ does not come from the Gilkey formula \emph{per se},
but from the \textbf{combination} of:
\begin{enumerate}
\item The Gilkey gauge-kinetic coefficient ($\propto q^2$);
\item The finite-space trace over SM representations;
\item Normalization conventions ($\alpha$ vs.\ $\kappa$, etc.).
\end{enumerate}

The question is whether these combine to give exactly $35/18$.

\subsection{Checking the Combination}

In the Chamseddine--Connes framework (with a specific spectral triple),
the gauge coupling for U(1) hypercharge is:
\begin{equation}\label{eq:CC-gauge}
\frac{1}{g_1^2} = \frac{f(0)}{2\pi^2} \times c_Y \times
a_4^{\rm gauge}(\mathcal{O}(1))
\end{equation}
where $c_Y = \mathrm{Tr}_F(Y^2/4)$ is the hypercharge trace.

For the Standard Model spectral triple:
\begin{equation}
c_Y = \frac{1}{4}\mathrm{Tr}_F(Y^2) \times g_F = \frac{10}{4} \times 8
= 20
\end{equation}
(But this depends sensitively on the grading conventions.)

With $a_4^{\rm gauge}(\mathcal{O}(1)) = 1/96$:
\begin{equation}
\frac{1}{g_1^2} = \frac{f(0)}{2\pi^2} \times 20 \times \frac{1}{96}
= \frac{f(0)}{2\pi^2} \times \frac{5}{24}
\end{equation}

Since $\beta_{U(1)} = 1/g_1^2$ and $\alpha^{-1} = 6\pi\beta_{U(1)}$
(using the \emph{approximate} $\kappa \approx \beta$ replacement,
which is already wrong at leading order), we get:
\begin{equation}
\alpha^{-1} = 6\pi \times \frac{f(0)}{2\pi^2} \times \frac{5}{24}
= \frac{6\pi \cdot 5 \cdot f(0)}{48\pi^2}
= \frac{5f(0)}{8\pi}
\end{equation}

This does not produce $35\pi/3$ or anything resembling 137.

The reason is clear: \textbf{the Chamseddine--Connes spectral action
requires the cutoff function $f$}, whose moments $f(0), f_2, f_4$
are free parameters.  The gauge coupling in that framework is
$1/g^2 \propto f(0)$, and $f(0)$ is \emph{not} determined by the
geometry.  In DFD, $f(0)$ is effectively replaced by the Chern--Simons
weighted average $\beta_{U(1)}$, but this replacement is
\textbf{not part of the Gilkey formula}.

%=============================================================================
\section{Where Does 35/18 Actually Come From?}
\label{sec:origin}
%=============================================================================

\subsection{Reconstruction of the Claimed Derivation}

Reading the existing documents carefully, the ``derivation'' of 35/18
proceeds as follows:

\begin{enumerate}
\item \textbf{Start with the target:} $\alpha^{-1} = 137.036$, so
  $\kappa_{U(1)} = 137.036/(6\pi) = 7.2700$.

\item \textbf{Divide by $\beta_{U(1)}$:}
  $\kappa/\beta = 7.2700/3.7969 = 1.91469$.

\item \textbf{Note that the Toeplitz factors simplify:}
  $(63/64)(4096/4095) = 64/65 = 0.98462$.

\item \textbf{Divide:} $1.91469 / 0.98462 = 1.94461$.

\item \textbf{Find the closest simple fraction:} $35/18 = 1.94\overline{4}$
  matches to 85 ppm.

\item \textbf{Assign physical meanings a posteriori:}
  $35 = 7 \times 5$, $18 = 2 \times 9$.
\end{enumerate}

This is \textbf{reverse engineering}, not derivation.  The fraction
$35/18$ was found by numerical search and then decomposed into
physically meaningful factors.

\subsection{The ``Physical Meanings'' Are Post-Hoc}

The decomposition $35/18 = (2/3)(5/2)(7/6)$ uses:
\begin{itemize}
\item $2/3$: ratio of $4\pi$ to $6\pi$ (normalization).
\item $5/2 = 1 + R_W/4$ where $R_W = 6$ is the Wilson ratio.
  But the Wilson ratio is $\beta_{SU(2)}/\beta_{U(1)}$, which is
  \emph{already used} in the electroweak mixing formula
  $\alpha^{-1} = 6\pi\kappa_{U(1)}$.  Using it again in the
  prefactor double-counts the electroweak mixing.
\item $7/6 = N_{\rm sp}/(N_{\rm gen} \cdot n_2)$: this is an SM
  quantum number ratio with no clear origin in the $a_4$ formula.
\end{itemize}

\subsection{The Alternative Decomposition $7 \times 5/(2 \times 9)$}

\begin{itemize}
\item $7 = N_{\rm sp}$: Standard Model SU(2) species count.
  This is a trace factor from the spectral triple, not from Gilkey.
\item $5 = n$: number of trivial bundles.  As shown in Section~5,
  these contribute \textbf{zero} to $a_4^{\rm gauge}$.
\item $2$: unclear origin (sometimes claimed as ``the 2 in 2a'',
  sometimes as normalization).
\item $9 = a = q_1^2$: the twist degree.  This enters the Gilkey
  formula as $|\Omega|^2 \propto q^2 = 81$, not as $q^2/q^2 = 1$.
  In $35/18$, the factor of 9 appears in the \emph{denominator},
  which means it is \emph{dividing out} the $q^2$ dependence.
  But $a_4^{\rm gauge} \propto q^2$, so dividing by $q^2$ removes
  the geometric information that the Gilkey formula provides.
\end{itemize}

%=============================================================================
\section{Quantitative Cross-Check}
\label{sec:crosscheck}
%=============================================================================

\subsection{What the Gilkey $a_4$ Gauge Term Gives Directly}

The gauge-kinetic coefficient for $\mathcal{O}(q)$ on
$\mathbb{C}P^2 \times S^3$ is:
\begin{equation}
\kappa_{\rm Gilkey}(q)
= \frac{q^2}{96} \times \frac{\sqrt{\pi}}{4}
= \frac{q^2\sqrt{\pi}}{384}
\end{equation}

For $q = 9$: $\kappa_{\rm Gilkey}(9) = 81\sqrt{\pi}/384 = 0.3738$.

\subsection{Ratio to Target}

The target is $\kappa_{U(1)} = 7.2700$.  The ratio:
\begin{equation}
\frac{\kappa_{U(1)}}{\kappa_{\rm Gilkey}(9)}
= \frac{7.2700}{0.3738} = 19.45
\end{equation}

This factor of $\sim 19.5$ must come from:
\begin{itemize}
\item The cutoff function normalization $f(0)/f_{\rm natural}$.
\item The Chern--Simons renormalization $\beta_{U(1)} = 3.797$.
\item The SM trace factors.
\item The Toeplitz corrections.
\end{itemize}

None of these are part of the Gilkey $a_4$ formula.  The Gilkey
formula gives $27\sqrt{\pi}/128$, and everything else is external
input.

%=============================================================================
\section{The Honest Assessment}
\label{sec:honest}
%=============================================================================

\begin{tcolorbox}[colback=yellow!10!white, colframe=orange!60!black,
  title=\textbf{Verdict: 35/18 Is Not Derived from the Gilkey $a_4$ Formula}]

\textbf{What the Gilkey $a_4$ formula provides:}
\begin{itemize}
\item The gauge-kinetic coefficient for $\mathcal{O}(q)$ on
  $\mathbb{C}P^2$: $a_4^{\rm gauge} = q^2/96$.
\item For $\mathcal{O}(9)$: $a_4^{\rm gauge} = 27/32$.
\item The five trivial bundles $\mathcal{O}^{\oplus 5}$ contribute zero.
\item On $\mathbb{C}P^2 \times S^3$: multiply by $a_0(S^3) = \sqrt{\pi}/4$.
\end{itemize}

\textbf{What the Gilkey formula does NOT provide:}
\begin{itemize}
\item The factor $N_{\rm sp} = 7$ (this is an SM trace over the
  spectral triple's finite Hilbert space).
\item The factor $n = 5$ in the \emph{numerator} of $35/18$ (the
  trivial bundles have $\Omega = 0$ and cannot contribute to the
  gauge coupling through $a_4^{\rm gauge}$).
\item The Chern--Simons renormalization $\beta_{U(1)}$.
\item The Toeplitz truncation corrections.
\item The electroweak mixing formula $\alpha^{-1} = 6\pi\kappa$.
\end{itemize}

\textbf{The factor 35/18 is a numerical approximation} that matches
$\kappa/[\beta \times 64/65]$ to 85 ppm.  Its decomposition into
$7 \times 5/(2 \times 9)$ uses SM quantum numbers that are external
to the Gilkey formula.  In particular, the appearance of $n = 5$
(trivial bundles) in the numerator is \textbf{inconsistent} with
the Gilkey computation, since trivial bundles do not contribute to
the gauge-kinetic $a_4$ coefficient.
\end{tcolorbox}

%=============================================================================
\section{What Would a Genuine Derivation Look Like?}
\label{sec:genuine}
%=============================================================================

A genuine derivation of the prefactor in $\kappa/\beta$ would need to:

\begin{enumerate}
\item \textbf{Specify the spectral triple} $(A, H, D)$ completely:
  the algebra $A$, the Hilbert space $H = L^2(M, S) \otimes H_F$,
  and the Dirac operator $D = D_M \otimes 1 + \gamma_5 \otimes D_F$.

\item \textbf{Compute the spectral action} $\mathrm{Tr}\,f(D^2/\Lambda^2)$
  using the heat-kernel expansion on the full product geometry
  $M^4 \times F$ (where $F$ is the finite noncommutative space).

\item \textbf{Extract the gauge coupling} from the $a_4$ coefficient
  of the full operator $D^2$, which involves both the gravitational
  and internal curvatures.

\item \textbf{Identify the cutoff} $\Lambda$ with the Toeplitz
  truncation scale, and show that this produces a specific numerical
  value for the gauge coupling.

\item \textbf{Show that $\beta_{U(1)}$} (the Chern--Simons weighted
  average) enters through the $S^3$ factor of the internal space
  in a specific, derivable way.
\end{enumerate}

The existing DFD documents do not perform steps (1)--(5) in detail.
Instead, the number $\kappa_{U(1)} = 7.2700$ is stated as the output
of a computation whose intermediate steps are not shown.  The
factorization into $35/18 \times (1+\Delta) \times \beta \times 64/65$
is a \emph{parametrization} of the result, not a derivation.

%=============================================================================
\section{Salvageable Content}
\label{sec:salvage}
%=============================================================================

Despite the above critique, several elements of the DFD framework
\emph{are} well-founded:

\begin{enumerate}
\item \textbf{$\alpha^{-1} = 6\pi\kappa_{U(1)}$}: This follows from
  standard electroweak mixing \emph{if} the Frame Stiffness Theorem
  $\kappa_{SU(2)} = 2\kappa_{U(1)}$ holds.  This relation is a
  property of the lattice action with the correct Wilson normalization.

\item \textbf{The Toeplitz factors $63/64$ and $4096/4095$}: These have
  clear algebraic origins in the matrix truncation of the algebra of
  functions on $\mathbb{C}P^2$ and the choice of regular vs.\ fundamental
  representation for the internal Hilbert space.

\item \textbf{$\beta_{U(1)} = 3.7969\ldots$}: This is a well-defined
  finite sum from the Chern--Simons partition function with $k_{\max} = 60$.

\item \textbf{$k_{\max} = 60$}: This follows from the Hirzebruch--Riemann--Roch
  theorem (Bridge Lemma): $\chi(\mathcal{O}(9) \oplus \mathcal{O}^{\oplus 5})
  = \binom{11}{2} + 5 = 60$.

\item \textbf{The $a_4$ Gilkey formula itself}: The formula is standard
  and correct.  The issue is not the formula but the claim that 35/18
  comes from it.
\end{enumerate}

%=============================================================================
\section{Summary}
\label{sec:summary}
%=============================================================================

\begin{table}[h]
\centering
\renewcommand{\arraystretch}{1.4}
\caption{Status of each factor in $\kappa/\beta \approx (35/18)(64/65)$.}
\begin{tabular}{llp{6cm}}
\toprule
\textbf{Factor} & \textbf{Value} & \textbf{Status} \\
\midrule
$64/65 = d/(d+1)$ & $0.98462$ &
\textcolor{green!60!black}{\textbf{Derived.}}
Traceless projection $\times$ regular-module boost, from Toeplitz truncation
at $d = 64$. \\
$\beta_{U(1)}$ & $3.7969$ &
\textcolor{green!60!black}{\textbf{Derived.}}
Chern--Simons weighted average with $k_{\max} = 60$. \\
$6\pi$ & $18.850$ &
\textcolor{green!60!black}{\textbf{Derived.}}
Electroweak mixing $+$ stiffness ratio. \\
$35/18$ & $1.944\overline{4}$ &
\textcolor{red!70!black}{\textbf{Not derived from Gilkey.}}
Reverse-engineered to fit $\kappa/(\beta \cdot 64/65) = 1.9446$.
The decomposition into $7 \times 5/(2 \times 9)$ assigns SM quantum
numbers to a numerically fitted fraction. \\
$1 + \Delta(d)$ & $1.000085$ &
\textcolor{orange!70!black}{\textbf{Empirical correction.}}
The coefficient $A = 0.3495$ in $\Delta = A/d^2$ is stated without
derivation. \\
\bottomrule
\end{tabular}
\end{table}

\begin{tcolorbox}[colback=blue!5!white, colframe=blue!50!black,
  title=\textbf{Final Conclusion}]

The prefactor $35/18 = 7 \times 5/(2 \times 9)$ is \textbf{not derived}
from the Gilkey $a_4$ formula for the connection Laplacian on
$\mathcal{O}(9) \oplus \mathcal{O}^{\oplus 5}$ over $\mathbb{C}P^2$.

\medskip\noindent
The Gilkey gauge-kinetic coefficient for this bundle is $27/32$, and
no algebraic manipulation of the Gilkey output (with or without the
$S^3$ factor, with or without normalization by $a_0$) produces $35/18$.

\medskip\noindent
The fundamental problem is that $n = 5$ (trivial bundles) appears in
the numerator of $35/18$, but trivial bundles have flat connections
and contribute \textbf{zero} to the gauge-kinetic $a_4$ coefficient.

\medskip\noindent
The factor $35/18$ is a \textbf{85 ppm numerical approximation} to
the ratio $\kappa/[\beta \times 64/65]$, found by searching for simple
fractions near $1.9447$.  Its decomposition into SM quantum numbers
is suggestive but \textbf{a posteriori}.

\medskip\noindent
This does not invalidate the DFD prediction $\alpha^{-1} = 137.036$,
which may still be correct through a computation that has not been
fully exhibited.  But the \emph{specific claim} that 35/18 ``comes from
the $a_4$ Gilkey gauge prefactor'' is \textbf{false}.
\end{tcolorbox}

%=============================================================================
\section*{References}
%=============================================================================

\begin{enumerate}
\item[{[1]}] P.~B.~Gilkey, \textit{Invariance Theory, the Heat Equation,
  and the Atiyah--Singer Index Theorem}, 2nd ed.\ (CRC Press, 1995).

\item[{[2]}] D.~V.~Vassilevich, ``Heat kernel expansion: user's manual,''
  \textit{Phys.\ Rep.}\ \textbf{388} (2003) 279--360,
  \href{https://arxiv.org/abs/hep-th/0306138}{arXiv:hep-th/0306138}.

\item[{[3]}] A.~H.~Chamseddine, A.~Connes, ``The spectral action principle,''
  \textit{Commun.\ Math.\ Phys.}\ \textbf{186} (1997) 731--750.

\item[{[4]}] G.~Alcock, ``Ab Initio Derivation of the Fine Structure Constant
  from Density Field Dynamics'' (2025).

\item[{[5]}] G.~Alcock, ``Density Field Dynamics: A Complete Unified Theory,''
  v108.
\end{enumerate}

\end{document}
